\begin{lemma}
Let \(C(x)=\binom{x}{2}_{\mathbb R}=x(x-1)/2\).  Assume
\[
0\le\varepsilon_1,\qquad 0\le\varepsilon_2\le1,\qquad
3\varepsilon_2\le\varepsilon_1/10,
\]
and let \(D,b,T\in\mathbb R\) satisfy
\[
D\ge98,\qquad b\ge0,\qquad T\ge0,\qquad
T\le(1+\varepsilon_2)D.
\]
Assume also that \(D\) and \(T\) are nonnegative integers, viewed as real
numbers.
Then
\[
C(T)\le(1+\varepsilon_1)C(D).
\]
Moreover, with \(M=(1+\varepsilon_2)b\), if \(T\le M\), then
\[
C(T)\le(1+\varepsilon_1)\bigl(C(b)+C(D-b)\bigr),
\]
and if \(M\le T\), then
\[
C(M)+C(T-M)
\le
(1+\varepsilon_1)\bigl(C(b)+C(D-b)\bigr).
\]
\end{lemma}
\begin{proof}
First record two elementary facts about \(C\).  For \(1\le u\le v\),
\[
C(v)-C(u)=\frac{(v-u)(u+v-1)}2\ge0,
\]
so \(C\) is monotone on \([1,\infty)\), and this is the only monotonicity
used below.  Also
\[
C(b)+C(D-b)
=\left(b-\frac D2\right)^2+\frac{D(D-2)}4. \tag{1}
\]
Since \(D\ge98\), the right hand side is nonnegative for every \(b\).

The hierarchy-type hypothesis gives the small numerical conversion
\[
10\varepsilon_2\le\varepsilon_1. \tag{2}
\]
Indeed \(3\varepsilon_2\le\varepsilon_1/10\) implies
\(30\varepsilon_2\le\varepsilon_1\).

For the one-block estimate, if \(T\le1\), then \(C(T)\le0\), while
\((1+\varepsilon_1)C(D)\ge0\) by \(D\ge98\) and \(\varepsilon_1\ge0\).  If
\(1\le T\), monotonicity on \([1,\infty)\) and
\(T\le(1+\varepsilon_2)D\) give
\[
C(T)\le C((1+\varepsilon_2)D).
\]
Expanding and using \(D\ge2\), \(0\le\varepsilon_2\le1\), gives
\[
\frac{C((1+\varepsilon_2)D)}{C(D)}
=(1+\varepsilon_2)
  \frac{(1+\varepsilon_2)D-1}{D-1}
\le 1+10\varepsilon_2.
\]
By (2), \(C((1+\varepsilon_2)D)\le(1+\varepsilon_1)C(D)\).  This proves the
one-block estimate.

We next prove the small-mass conditional estimate.  Write
\[
Q(b)=C(b)+C(D-b).
\]
By (1), \(Q(b)\ge0\).  If \(T\le1\), then \(C(T)\le0\), so the desired
inequality follows.  Assume \(T\ge1\) and \(T\le M=(1+\varepsilon_2)b\).
There are two regimes.

First suppose \(D-b\le0\) or \(1\le D-b\).  The following
cleared-denominator estimate applies:
\[
2\bigl((1+10\varepsilon_2)Q(b)-C((1+\varepsilon_2)b)\bigr)
 =
(D-b)(D-b-1)+
\varepsilon_2\bigl(18b^2-20Db+10D^2-10D+b\bigr)
-\varepsilon_2^2b^2. \tag{3}
\]
On this regime the first term is nonnegative.  The remaining
\(\varepsilon_2\)-part is also nonnegative: since \(\varepsilon_2\le1\), it
is enough to check
\[
17b^2-20Db+10D^2-10D+b\ge0.
\]
This quadratic in \(b\) has positive leading coefficient and discriminant
\[
(-20D+1)^2-68(10D^2-10D)=-280D^2+640D+1,
\]
which is negative for \(D\ge98\).  Thus
\[
C((1+\varepsilon_2)b)
\le(1+10\varepsilon_2)Q(b). \tag{4}
\]
Monotonicity on \([1,\infty)\), applied only because \(1\le T\le M\), gives
\[
C(T)\le C(M)=C((1+\varepsilon_2)b)\le(1+10\varepsilon_2)Q(b).
\]
Using (2) proves the desired bound in this first regime.

It remains to treat the only regime where a real binomial term can become
negative:
\[
0<D-b<1. \tag{5}
\]
Write \(D=d\) and \(T=t\) with \(d,t\in\mathbb N\), and put
\(u=d-b\), so \(0<u<1\).  If \(t\le d-1\), then monotonicity on
\([1,\infty)\) and the identity
\[
Q(d-u)-C(d-1)=(1-u)(d-u-1)\ge0,
\qquad u=d-b\in(0,1),
\]
give \(C(T)\le Q(b)\), hence the result.  It remains to consider all
integer values \(t\ge d\).  Write \(t=d+r\), where \(r\in\mathbb N\).  The
small-mass hypothesis \(T\le(1+\varepsilon_2)b\) becomes
\[
d+r\le(1+\varepsilon_2)(d-u),
\]
or equivalently
\[
r+u\le\varepsilon_2(d-u)\le\varepsilon_2d. \tag{6}
\]
Thus \(\varepsilon_2\ge(r+u)/d\), and \(r\le d\).  Moreover
\[
Q(d-u)-C(d-1)=(1-u)(d-u-1)\ge0,
\]
so \(Q(b)\ge C(d-1)\).  Finally,
\[
\begin{aligned}
C(d+r)-Q(d-u)
&= rd+\frac{r(r-1)}2+u(d-u)  \\
&\le \frac32\,d(r+u).
\end{aligned}
\]
Combining this with (6) and \(Q(b)\ge C(d-1)\), and using \(d\ge98\), gives
\[
C(T)-Q(b)
\le \frac32\,d(r+u)
\le \frac32\,\varepsilon_2 d^2
\le 30\varepsilon_2 C(d-1)
\le 30\varepsilon_2 Q(b)
\le \varepsilon_1 Q(b),
\]
where the last inequality is (2).  Hence
\(C(T)\le(1+\varepsilon_1)Q(b)\).  This covers every allowed integer
\(T=t\) in the exceptional interval (5), including \(t\ge d+1\), and
therefore completes the small-mass conditional estimate.

It remains to prove the large-mass conditional estimate.  Suppose \(M\le T\).
Together with \(T\le(1+\varepsilon_2)D\), this implies \(b\le D\), and hence
\[
0\le T-M\le(1+\varepsilon_2)(D-b). \tag{7}
\]
Put \(u=D-b\) and \(y=T-M\).  If \(1\le y\), then monotonicity on
\([1,\infty)\) and (7) give
\[
C(y)\le C((1+\varepsilon_2)u).
\]
The first block is exactly \(C(M)=C((1+\varepsilon_2)b)\).  For the two
stretched blocks together, clearing denominators gives
\[
\begin{aligned}
&2\Bigl((1+10\varepsilon_2)\bigl(C(b)+C(D-b)\bigr)
  -C((1+\varepsilon_2)b)
  -C((1+\varepsilon_2)(D-b))\Bigr)  \\
&\qquad =
\varepsilon_2\Bigl((8-\varepsilon_2)
  \bigl(b^2+(D-b)^2\bigr)-9D\Bigr).
\end{aligned}
\]
Because \(0\le b\le D\), the sum \(b^2+(D-b)^2\) is at least \(D^2/2\).
Together with \(D\ge98\) and \(\varepsilon_2\le1\), the right hand side is
nonnegative.  Hence, using (2),
\[
C((1+\varepsilon_2)b)+C((1+\varepsilon_2)(D-b))
\le
(1+\varepsilon_1)\bigl(C(b)+C(D-b)\bigr).
\]
Combining the preceding two displays proves the large-mass conditional
estimate when \(1\le y\).

It remains to handle \(0\le y\le1\).  If \(u=0\) or \(1\le u\), then
\(C(y)\le0\), and estimate (4), applied to
\(C((1+\varepsilon_2)b)\), gives
\[
C(M)+C(y)\le(1+\varepsilon_1)Q(b).
\]
The only leftover case is
\[
0<u=D-b<1,\qquad 0\le y\le1. \tag{8}
\]
Write again \(D=d\).  Since \(M=(1+\varepsilon_2)(d-u)\ge d-u>d-1\) and
\(M\le T\), the integer \(T=t\) satisfies \(t\ge d\).  Write
\(T=t=d+r\) with \(r\in\mathbb N\).  Since
\(T\le(1+\varepsilon_2)D\), we have
\[
r\le\varepsilon_2 d. \tag{9}
\]
Also \(y-u=r-\varepsilon_2b\), so
\[
|y-u|\le r+\varepsilon_2b\le2\varepsilon_2d. \tag{10}
\]
For \(0\le a,c\le1\),
\[
C(a)-C(c)=\frac{(a-c)(a+c-1)}2\le\frac{|a-c|}{2}.
\]
Using this with \(a=y\) and \(c=u\), and then using (10), gives
\[
C(y)-C(u)\le\varepsilon_2d. \tag{11}
\]
For the first block,
\[
C((1+\varepsilon_2)b)-C(b)
=\frac{\varepsilon_2b((2+\varepsilon_2)b-1)}2
\le\frac32\,\varepsilon_2d^2, \tag{12}
\]
because \(0\le b\le d\) and \(\varepsilon_2\le1\).  Adding (11) and (12),
and using \(d\ge98\), gives
\[
C(M)+C(y)-Q(b)\le 2\varepsilon_2d^2
\le 30\varepsilon_2 C(d-1)
\le 30\varepsilon_2 Q(b)
\le \varepsilon_1 Q(b).
\]
This proves \(C(M)+C(T-M)\le(1+\varepsilon_1)Q(b)\) in the exceptional
large-mass subcase (8), and completes the proof.
\end{proof}
